\begin{apartats}

\apartat[1,25]{1}
Estudia la curvatura i troba els punts d'inflexió de la funció
\[
f(x)=x^3-6x^2+9x+1 .
\]

\begin{solucio}
$f'(x)=3x^2-12x+9$ i $f''(x)=6x-12$, que s'anul·la en $x=2$.\\
$f''<0$ a $(-\infty,2)$: \textbf{còncava} (cap avall). \quad $f''>0$ a $(2,+\infty)$:
\textbf{convexa} (cap amunt).\\
Com que la curvatura hi canvia i $f(2)=8-24+18+1=3$, el punt d'inflexió és $\boxed{(2,3)}$.
\end{solucio}

\begin{nomesllarg}
\apartat{0,75}
Estudia la curvatura de la funció $g(x)=\ln\left(x^2+1\right)$ i troba'n els punts d'inflexió.

\begin{solucio}
$g'(x)=\dfrac{2x}{x^2+1}$ i
$g''(x)=\dfrac{2\left(x^2+1\right)-2x\cdot2x}{\left(x^2+1\right)^2}
=\dfrac{2-2x^2}{\left(x^2+1\right)^2}$.\\
El denominador és positiu, així que el signe és el de $2-2x^2$: \textbf{convexa} a $(-1,1)$ i
\textbf{còncava} fora.\\
Punts d'inflexió: $(-1,\ln2)$ i $(1,\ln2)$.
\end{solucio}
\end{nomesllarg}

\apartat[1,25]{0,75}
Pot una funció polinòmica de grau $3$ no tenir cap punt d'inflexió? Raona-ho.

\begin{solucio}
Si $f(x)=ax^3+bx^2+cx+d$ amb $a\neq0$, aleshores $f''(x)=6ax+2b$, que és una recta de pendent
$6a\neq0$.\\
Una recta no horitzontal s'anul·la exactament un cop i hi canvia de signe. Per tant, una
funció polinòmica de grau $3$ té \textbf{sempre exactament un punt d'inflexió}, i no pot no
tenir-ne cap.
\end{solucio}

\end{apartats}
